\chapter*{\textbf{Resumen}}
\vspace{-2cm}

\begingroup
\setlength{\parskip}{0.8em}
\setlength{\parindent}{0pt}

Este trabajo presenta la implementación de un sistema de narración automática de partidos de fútbol a partir de vídeo broadcast. El sistema se organiza como un pipeline que combina visión por computador, tracking multiobjeto, análisis geométrico, detección de eventos y generación de comentarios en texto y audio. Su objetivo es extraer información suficiente para identificar qué acción ocurre durante el partido, qué jugador participa en ella y en qué contexto futbolístico se produce.

La principal innovación de este trabajo no es solo cada módulo por separado, sino haberlos integrado en un pipeline completo capaz de convertir un vídeo de fútbol en una narración automática. Es algo completamente novedoso que no se había implementado en ningún trabajo previamente.

La primera fase del sistema se centra en el procesamiento secuencial del video frame a frame. En esta etapa se realiza la detección y seguimiento de los elementos relevantes del juego, principalmente jugadores, árbitros y balón. Para ello se utiliza un detector basado en YOLO y un módulo de tracking que combina ByteTrack con una capa de tracking canónico, lo que permite mantener identidades consistentes a lo largo del vídeo. Además, se realiza una identificación de equipo mediante el color de la camiseta, estimación de la posición del jugador sobre el terreno de juego mediante homografía, apoyada en PnLCalib, y una clasificación posicional basada en un modelo de Set Transformer entrenado especificamente para aproximar roles de jugadores como lateral, mediocentro o delantero.

A partir de las trayectorias y estados generados, comienza una segunda fase de procesamiento asíncrono orientada a analizar la evolución espacial y temporal de los jugadores y del balón. En esta etapa se infieren acciones futbolísticas como pases, controles o tiros mediante un detector de acciones basado en PathCRF. Los eventos detectados sirven como entrada para un módulo de generación lingüística basado en un LLM cuantizado, encargado de producir comentarios adaptados al contexto del partido.

Finalmente, los comentarios generados se convierten a audio mediante síntesis de voz, usando sistemas como XTTS o ElevenLabs. Con el fin de producir una narración más natural y dinámica, el sistema contempla dos voces hispanohablantes, masculina y femenina, y combina comentarios sobre acciones del partido con información contextual y datos de la competición.

El pipeline se ha desarrollado con una orientación clara hacia la ejecución local en GPU, ya que la latencia constituye un requisito fundamental para aproximarse a un escenario de narración en tiempo real. Por ello, las distintas etapas del sistema buscan equilibrar precisión en el análisis del juego, calidad narrativa y coste computacional, tomando como referencia un entorno equipado con una GPU RTX 3080 de 10 GB de VRAM.

El pipeline se expone mediante dos vías de consumo: una herramienta de línea de comandos (track.py) orientada a procesamiento por lotes y experimentación, y una interfaz web que permite cargar vídeos, configurar alineaciones y obtener el vídeo narrado con el audio integrado, facilitando la demostración interactiva del sistema.

Los resultados muestran que la integración de detección, tracking, análisis geométrico y generación de lenguaje permite construir una narración automática y satisfactoria teniendo en cuenta las limitaciones existentes.

\section*{Palabras clave}

\noindent Narración automática, visión por computador, fútbol, detección, tracking multiobjeto, homografía, clasificación de acciones, Set Transformer, YOLO, generación de comentarios, síntesis de voz.

   
